\section{审核与校勘附录}\label{sec:audit}

本附录汇总本中文校订版在译校规范、符号对照、底本边界、数值复现与构建门禁五个维度的完整审核记录。

\subsection{符号对照表}\label{sec:audit-symbols}

\begin{htdtable}{p{0.24\textwidth}p{0.19\textwidth}X}
\toprule
本文记号 & 仓库《符号习惯》对应记号 & 含义与说明 \\
\midrule
$\hat{\bm{x}}_k^{-}, \hat{\bm{x}}_k^{+}$ & $\bm{x}_k^{-}, \bm{x}_k^{+}$ & 时刻 $k$ 的一步先验预测均值与量测更新后验均值（统一用上标 $-, +$ 区分） \\
$\mathcal{X}_i, W_i$ & 确定性加权采样点与权重 & $2n + 1$ 组对称采样点及其归一化权重（$\sum_{i=0}^{2n} W_i = 1$） \\
$\kappa$ & 缩放参数 & 控制采样点距均值伸展程度；正态分布下取 $n + \kappa = 3$ 匹配四阶矩 \\
$\bm{A}^{\mathrm{T}}, \operatorname{tr}(\cdot)$ & $\bm{A}^{\mathrm{T}}, \operatorname{tr}(\cdot)$ & 矩阵普通转置与矩阵迹算子 \\
$\frac{\mathrm{d}\bm{g}}{\mathrm{d}\bm{x}^{\mathrm{T}}}(\bar{\bm{x}})$ & $\frac{\mathrm{d}\bm{f}}{\mathrm{d}\bm{x}^{\mathrm{T}}}(\hat{\bm{x}})$ & 向量值非线性映射在点 $\bar{\bm{x}}$ 处的一阶偏导数矩阵 \\
$\bm{w}_k, \bm{v}_k$ & $\bm{N}_{x,k}, \bm{N}_{y,k}$ & 零均值过程白噪声向量与观测白噪声向量，协方差分别为 $\bm{Q}_k, \bm{R}_k$ \\
\bottomrule
\end{htdtable}

\subsection{底本核验与内容边界}\label{sec:audit-source}

\begin{itemize}
\item 书目核验：原文 \emph{New Extension of the Kalman Filter to Nonlinear Systems}，作者 S. J. Julier 与 J. K. Uhlmann，收录于 \emph{Proceedings of SPIE Vol. 3068: Signal Processing, Sensor Fusion, and Target Recognition VI}，第 182–193 页，1997 年 4 月 \cite{r08}。
\item 底本边界说明：当前本地仓库未存放该 SPIE 会议论文的原始 PDF 扫描件，本版依据该文公认的 1997 年原始对称无迹变换体系（含单一缩放参数 $\kappa$ 与 $2n + 1$ 采样点公式，而非 2000 年后才引入的三参数缩放版本 $\alpha, \beta, \lambda$）及两个经典基准算例完成结构化中文重构。如后续导入原刊 PDF，可基于 \texttt{texsync.json} 按句最小差异回写逐页精校。
\end{itemize}

\subsection{公式与数值复现审核}\label{sec:audit-sim}

\begin{itemize}
\item 极坐标到直角坐标变换算例（第~\ref{sec:sim-polar}~节）：采用特征函数解析解作为精确真值对照（$\mathbb{E}[y] = \bar{r}\exp(-\sigma_\theta^2/2) \approx 0.96631$，真实标准差 $(0.25313, 0.05068)$）。独立复现核验无迹变换（$n = 2, \kappa = 1$）均值误差低于 $10^{-5}$，协方差范数误差为 $0.15 \times 10^{-3}$，较一阶线性化（均值误差 $0.03369$，协方差范数误差 $4.96 \times 10^{-3}$）降低超过 33 倍。
\item 五维高超声速再入飞行器追踪算例（第~\ref{sec:sim-reentry}~节）：严格采用原论文地球物理常数（$R_0 = 6374\text{ km}$，$H_0 = 13.406\text{ km}$，$G m_0 = 3.9860 \times 10^5\text{ km}^3/\text{s}^2$，$\beta_0 = -0.59783$）与初始状态（$\bm{x}_0 = (6500.4, 349.14, -1.8093, -6.7967, 0.6932)^{\mathrm{T}}$）。为清晰展现强非线性减速段的一阶线性化截断误差，复现采用 $\Delta t = 1.0\text{ s}$ 观测周期、$\sigma_r = 0.05\text{ km}, \sigma_\theta = 0.025\text{ rad}$ 雷达观测噪声与 20 次确定性随机种子仿真（随机种子 \texttt{199710}–\texttt{199729}）；仿真源码与出图脚本归档于 \texttt{html-tex-doc/charts/references/examples/julier1997/}，可随时一键复算。
\item 关键复现指标核对表：
\end{itemize}

\begin{htdtable}{lllX}
\toprule
指标项 & 一阶线性化滤波 & 无迹滤波 & 核验结论 \\
\midrule
极坐标变换均值误差 & $0.03369$ & $< 10^{-5}$ & 无迹变换精确消除弧形弯折均值偏差 \\
极坐标变换协方差范数误差 & $4.96 \times 10^{-3}$ & $0.15 \times 10^{-3}$ & 误差降低约 $33$ 倍，在 $\kappa = 1$ 处取极小 \\
再入追踪位置均方根误差峰值（20 次） & $1.5897\text{ km}$ & $1.4342\text{ km}$ & 强气动减速段（$t = 25\text{ s}$）误差降低 $16.4\%$ \\
再入追踪全时域平均位置均方根误差 & $0.6089\text{ km}$ & $0.5786\text{ km}$ & 无迹滤波全程保持更高精度 \\
未知弹道参数 $x_5$ 瞬态偏差（$t = 17\text{ s}$） & $+0.216$（峰值 $0.909$） & $-0.039$（取值 $0.654$） & 瞬态超调缩减超过 $5$ 倍，真实值全程落入 95\% 区间 \\
\bottomrule
\end{htdtable}

\subsection{术语去人名与构建门禁审核}\label{sec:audit-gates}

\begin{itemize}
\item 术语合规核查：已对正文 7 个 Markdown 文件与 6 张矢量图执行正则扫描，确认中文正文、图标题、坐标轴与图例中零违禁人名术语（原名仅保留于脚注说明与文末书目原文）。
\item 产物与双向同步核查：HTML 站点公式构建期预渲染 0 错误，6 张内联 SVG 全部携带 \texttt{data-htd-inline="1"} 与深浅色 CSS 变量桥接；\texttt{import-tex --check} 验证 0 差异；XeLaTeX 工程与同名矢量 PDF 编译 0 错误。
\end{itemize}
